\section{AutoMF: model selection and ensemble fitting}
\label{sec:method}
\subsection{From parameter inputs to field predictions}
At fidelity $\ell$, the training data are
$\calD_\ell=\{(x_i^{(\ell)},y_i^{(\ell)})\}_{i=1}^{N_\ell}$: $N_\ell$ examples
indexed by $i$, each pairing $d$ simulation parameters $x_i^{(\ell)}\in\R^d$
with a field $y_i^{(\ell)}$ on grid $\Omega_\ell$. The index $\ell$ identifies
fidelity. Each surrogate model $f_m$ maps the parameter vector $x$
to a fine field on a common target grid. For models that incorporate a coarse
field, this field is obtained internally: either a trained low fidelity
surrogate model predicts it at $x$, or nearest neighbour retrieval selects the stored
coarse solution whose training parameters are closest to $x$ in standardized
parameter space. The coarse field then serves as a reference for a learned
correction. Thus, $f_m$ denotes the complete inference procedure, including any
coarse prediction, retrieval, and correction, and requires no additional
simulation at the query input.

We reserve $K$ additional fine examples for \emph{ensemble fitting} and a separate evaluation set.
We refer to the reserved input and fine field pairs as \emph{fitting examples}.
Both sets are excluded from low fidelity (LF) and high fidelity (HF) training
before normalization, basis construction, or model fitting. Input identities
establish correspondence between fidelities. Each dataset has its own trained
surrogate library and ensemble weights. Each rule determines its weights from
the fitting examples and saves them before scoring evaluation fields.
Appendix~\ref{app:splits} discusses split integrity and prior access to
historical evaluation results.

\subsection{Surrogate model library}
The surrogate library contains nine models, M1 to M9 in Table~\ref{tab:paperbaselines},
chosen to vary both the representation of a field and the way coarse data
enter learning. Transfer can reuse features across fidelities, joint supervision
shares information during training, and correction uses a coarse field as a
spatial reference. A fine only reduced basis model supplies an alternative
when coarse information is less useful. We adapt these strategies to the same
parameter input and target field interface, including the internal coarse
prediction or retrieval required by correctors. Three models first learn from coarse fields
and then finetune on fine fields, following multifidelity transfer learning
\citep{lyu2023}. FiLM FNO transfer combines Fourier mixing
\citep{li2021fno} with parameter conditioned featurewise linear modulation
(FiLM) \citep{perez2018}.
ConvNeXt transfer uses convolutional blocks in an encoder and decoder with skip
connections \citep{liu2022convnext}. Wavelet transfer uses a custom wavelet
representation inspired by WNO \citep{tripura2022wno}. Together they test coarse
initialization across different spatial representations. All pairs FNO instead
pools supervision across fidelity pairs before fine training. It takes
parameters and two fidelity labels as inputs. Poseidon motivates pair enumeration
\citep{herde2024}, while our FNO \citep{li2021fno} uses neither its temporal
operator nor a supplied source field.

Four models use coarse predictions or retrieved fields to guide a correction.
Distribution FNO predicts a residual from coarse ensemble summaries, adapting
FIRE's distribution conditioning idea \citep{yu2026fire}. Retrieved field
DeepONet uses a branch and trunk representation \citep{lu2021,lu2022mf} to
correct a neighboring coarse training field. The slice attention and ConvNeXt
correctors refine a predicted coarse field through six updates, using modified
Transolver inspired attention \citep{wu2024transolver} and convolutional blocks
\citep{liu2022convnext}, respectively. Their coarse inputs are assembled out of
fold with four ensemble members per case. Finally, POD GP predicts coefficients
in a proper orthogonal decomposition (POD) basis using Gaussian processes (GP)
\citep{higdon2008}. It tests whether a compact field representation can be
learned directly from fine solutions. Our adaptations specify how parameters
condition each model, how fidelities supervise it, and how coarse references
are constructed. Appendix~\ref{app:adaptations} explains for each model
the source components, our changes, and the inclusion rationale. Architecture and training cost
both vary across the library, so rankings cannot be attributed to one design
choice alone.

\subsection{Fitting one weight per model}
The ensemble averages the model predictions,
\begin{equation}
 \widehat y_w(x)=\sum_{m=1}^M w_m f_m(x),\qquad
 w\in\simplex=\{w\in\R^M:w_m\geq0,\ \mathbf{1}^\top w=1\}.
 \label{eq:mixture}
\end{equation}
Here $M=9$ counts models, $w_m$ weights model $m$, and $\widehat y_w$ is the
combined field. The simplex $\Delta_M$ contains nonnegative weights summing to
one. The vector $\mathbf{1}$ contains ones and $\top$ denotes transpose.
Weights stay fixed across inputs and grid points. For fitting pair $(x_i,y_i)$,
define $\langle a,b\rangle_Q=a^\top Qb$, $\norm a_Q=\sqrt{a^\top Qa}$, and
$s_i=\max(\norm{y_i}_Q,10^{-8})$. Fields are vectors of grid values, $Q=I$ is
the identity matrix, and the floor prevents division by zero. We compute
\begin{equation}
 e_{im}=\frac{f_m(x_i)-y_i}{s_i},\qquad
 [G_i]_{mn}=\langle e_{im},e_{in}\rangle_Q,\qquad
 \widehat C=\frac{1}{K}\sum_{i=1}^K G_i.
 \label{eq:gram}
\end{equation}
Here $e_{im}$ is model $m$'s relative error field on example $i$, $G_i$ records
error products between models $m$ and $n$, and $\widehat C$ averages these over
$K$ fields. Diagonal entries measure individual squared error. Other entries
measure error alignment, retaining information needed to score a mixture.

\noindent\textbf{Selected model.} This rule gives weight one to model $m^*$
with the lowest mean relative $L^2$ error on the fitting examples,
$m^*=\arg\min_m K^{-1}\sum_i\sqrt{[G_i]_{mm}}$.

\noindent\textbf{Inverse error mixture.} This rule gives more weight to individually
accurate models:
\begin{equation}
 w_m^{\rm inv}=\frac{1/\max(\widehat C_{mm},\tau)}
 {\sum_j 1/\max(\widehat C_{jj},\tau)},\qquad \tau=10^{-30}.
\label{eq:inverse}
\end{equation}
Here $w_m^{\rm inv}$ is the normalized inverse error weight, $j$ runs over all
models, and $\tau$ prevents division by zero. The weights sum to one.

\noindent\textbf{Fitted mixture.} This rule chooses the weights together by minimizing
the error of their combined prediction:
\begin{equation}
 \widehat w^{\rm full}\in\arg\min_{w\in\simplex}w^\top\widehat Cw.
 \label{eq:stack}
\end{equation}
Here $\widehat w^{\rm full}$ minimizes the mean squared relative error
$w^\top\widehat Cw$ on the fitting examples. This is convex stacking
\citep{wolpert1992,breiman1996}, allowing a less accurate model to offset another's errors.
Both mixture rules use squared normalized errors, whereas model selection and
reporting use mean relative $L^2$. The objectives can have different optima.
We test the consequence of this choice in Section~\ref{sec:results} and derive
the distinction in Appendix~\ref{app:geometry}. Numerical details appear in
Appendix~\ref{app:numerics}.

Ensemble fitting estimates only nine weights, with eight free values because
they sum to one. The surrogate models remain fixed, and each fitting field
supplies spatial error products that summarize individual accuracy and error
alignment. Once these products are computed, the optimization uses a
$9\times9$ matrix regardless of the number of grid points. This small fitting
problem makes limited additional fine data a practical option. Spatial values
within a field remain dependent, however, so field resolution does not count
as additional independent fitting examples or guarantee generalization from
five fields. Appendix~\ref{app:algorithm} gives the pseudocode.
